\subsection{Critical Ricci control with a polar Jacobian bound}
\label{subsec:critical-radial-defect}

On the pre-cut polar coordinates from the physical center \(x_i\),
write \(t=d_{g_i}(x_i,\cdot)\), \(l=\log t\), and
\(p_i(l,\vartheta)=t^{-3}J_i(t,\vartheta)\), where \(J_i\) is
the polar volume Jacobian. Let \(\Theta_{i,l}\) be the unit
directions whose cut time is strictly beyond \(e^l\), and
extend \(p_i\) by zero after each cut. The following condition
uses the same directions at two radii:
\[
 Q_i(s,l)=\int_{\Theta_{i,l}}
         \frac{p_i(l,\vartheta)^2}{p_i(s,\vartheta)}
         \,d\vartheta,\qquad s<l .
\]

\begin{proposition}[A critical radial comparison]
\label{prop:critical-radial-defect}
Use the physical sequence and complete marks of
\eqref{eq:anchored-marks}. Suppose there are intervals
\([A_i,B_i]\subset[\log a_i,\log b_i]\) of lengths
\(T_i=B_i-A_i\to\infty\) and fixed \(C<\infty\) and
\(0\le\gamma<1\) such that
\begin{equation}\label{eq:critical-polar-jacobian}
 Q_i(s,l)\le C e^{2\gamma(l-s)}
 \qquad(A_i\le s<l\le B_i).
\end{equation}
There are scales \(a_i\ll\rho_i\ll b_i\) on which the
physical full curvature integral and total variation of
\(4dV_{h_i}-\Delta_{h_i}(d_{h_i}(x_i,\cdot)^2/2)\)
vanish on every fixed compact strip of positive radii,
with \(h_i=\rho_i^{-2}g_i\).
If the proper flat comparison with compact first-order
embeddings of Proposition~\ref{prop:radial-distance-test}
is supplied at these scales, its round cut has the two
complete marks on opposite sides and off the closed trim.
In the minimizing band this yields the contradiction of
that proposition.
\end{proposition}

\begin{proof}
On a surviving ray set
\(w=t\Delta_g t-3\) and
\(f=t^2(\Ric_g(\partial_t,\partial_t))_-\).
Riccati and the trace inequality for the three-dimensional
shape operator give \((\log p)'=w\) and
\(w'+w+w^2/3\le f\) in log radius. The critical volume
identity on the pre-cut region is
\[
 f^2p\,dl\,d\vartheta=(\Ric_g(\partial_t,\partial_t))_-^2dV_g .
\]
Integrate \(w'+w\le f\) backwards from \(l\), then
integrate the result over all start times in \([s,l]\). This gives
\[
 w(l)_+\le
 \frac{\log^+(p(l)/p(s))}{e^{l-s}-1}
 +\int_s^l\frac{e^{u-s}-1}{e^{l-s}-1}f(u)\,du .
\]
All radii in this inequality are before the given ray's cut.
Put \(P(l)=\int p(l)\) and \(G(l)=\int w(l)_+p(l)\), with
angular integration on surviving directions, and
\(F(u)=\int f(u)^2p(u)\). Jensen on the final angular
measure \(p(l)d\vartheta/P(l)\) and Cauchy--Schwarz with the
earlier weight \(p(u)\) give
\[
 G(l)\le
 \frac{P(l)\log(1+Q(s,l)/P(l))}{e^{l-s}-1}
 +\int_s^l
   \frac{e^{u-s}-1}{e^{l-s}-1}\sqrt{F(u)Q(u,l)}\,du .
\]
The first term is taken to be zero if \(P(l)=0\).
By Fatou as \(s\uparrow l\), \eqref{eq:critical-polar-jacobian}
also gives \(P(l)\le C\). For \(l-s\ge1\) the first term is
bounded by \(C'(1+l-s)e^{-(l-s)}\), and the kernel multiplying
\(\sqrt{F(u)}\) is bounded by \(C'e^{-(1-\gamma)(l-u)}\).

Restore the index and take \(s=A_i\). On the middle half
\(J_i\) of the log interval, integrate this estimate. Since
\(\int_{A_i}^{B_i}F_i\) is bounded by a constant times the
physical full curvature budget, Fubini and Cauchy--Schwarz
give \(\int_{J_i}G_i=O(1+\sqrt{T_i})\).
The normalized coarea function \(P_i\) has distributional
derivative \(\int w p_i\,d\vartheta\,dl-d\kappa_i\), where
\(\kappa_i\ge0\) records the downward jumps of the polar
density at cut times. This follows by differentiating
along each ray, with zero extension after cut, and then
integrating in angle. Thus its positive variation is
at most \(\int G_i\), and its total variation on \(J_i\)
is at most \(C+2\int_{J_i}G_i\), since \(0\le P_i\le C\).

Divide \(J_i\) into \(\lfloor T_i^{3/4}\rfloor\) equal
open subintervals. On one, the sum of the full curvature
integral, \(\int G_i\), and the variation of \(P_i\)
is \(O(T_i^{-1/4})\). Let \(\rho_i\) be the exponential
of its midpoint. Its width in log radius tends to
infinity, and the middle-half buffers separate the
scale from both quantiles. Bounded diameter gives
\(\rho_i\to0\). Put \(l_i=\log\rho_i\).

For each fixed \(0<a<b\), the physical rescaled volume
profile \(V_i(t)\) on \([a,b]\), up to an additive
constant, is \(\tfrac14 P_i(l_i)t^4+o(1)\), uniformly.
Indeed \(V_i'(t)=t^3 P_i(l_i+\log t)\) almost everywhere
and one can take a one-sided representative of \(P_i\).
The singular part of the distance Laplacian is nonpositive
\cite[Corollaries~2.11--2.12]{MantegazzaMascellaniUraltsev2014DistanceHessian}.
Thus the negative mass of the full rescaled defect
on \((a,b)\) is just
\(\int_{\log a}^{\log b}e^{4u}G_i(l_i+u)\,du=o(1)\).
The finite estimate of Lemma~\ref{lem:physical-flux}
gives the total variation claim on each inner trim.
The comparison tests for Proposition~\ref{prop:radial-distance-test}
lie in such trims and have bounded supremum norms.
Consequently their transport tests tend to zero, including
when the comparison images omit physical volume. That
proposition gives the marked cut and cap comparison.
\end{proof}

The centered and protected Ricci hypotheses of
Proposition~\ref{prop:physical-ricci-flux} imply
\eqref{eq:critical-polar-jacobian} with \(\gamma=0\):
polar Riccati bounds the forward ratio on surviving rays,
and in the protected case the measured transition bounds
the shifted area ratio. The weighted condition can
also hold in the presence of pointwise Ricci spikes.
At the critical exponent the usual polar multiplier
\(\psi^3 J\), for \(\psi=(\Delta_g t-3/t)_+\),
has no coercive \(\psi^4J\) term. Its curvature term
therefore lacks the estimate used in supercritical
comparison \cite[Theorem~2.1]{DaiWeiZhang2018Local}.
Condition~\eqref{eq:critical-polar-jacobian} instead
controls the cost of transporting the critical integral
from an earlier radius. It and the compact comparison
embeddings remain additional to the physical ball lower
bound and full curvature budget.
